% Instance identity: derived namespaces for co-resident instances (build/instance-identity).
\documentclass[11pt]{article}
\usepackage[T1]{fontenc}
\usepackage{lmodern}
\usepackage{microtype}
\usepackage[margin=1in]{geometry}
\usepackage{amsmath,amssymb,amsthm,mathtools}
\usepackage{algorithm}
\usepackage{algpseudocode}
\usepackage{booktabs}
\usepackage[numbers,sort&compress]{natbib}
\usepackage[hidelinks]{hyperref}
\usepackage[capitalise,noabbrev]{cleveref}

\theoremstyle{plain}
\newtheorem{theorem}{Theorem}[section]
\newtheorem{lemma}[theorem]{Lemma}
\newtheorem{proposition}[theorem]{Proposition}
\newtheorem{corollary}[theorem]{Corollary}
\theoremstyle{definition}
\newtheorem{definition}[theorem]{Definition}
\theoremstyle{remark}
\newtheorem{remark}[theorem]{Remark}

\newcommand{\lean}[1]{\texttt{#1}}
\newcommand{\env}{\texttt{MIND\_PALACE\_OVERLAY}}
\newcommand{\near}{\bowtie}
\newcommand{\shared}{\mathrm{shared}}
\newcommand{\loc}{\mathrm{loc}}

\title{Instance Identity: Disjoint Namespaces for Co-Resident Instances of a Knowledge Store}
\author{mind-palace contributors}
\date{7 November 2026}

\begin{document}
\maketitle

\begin{abstract}
A second instance of a knowledge store on one host, or a copied checkout, resolves every name in
its configuration (data root, tuned parameters, outbound lane, secret names, service labels, backup
repository, secrets server) to the first instance's. We require each named overlay to declare an
instance id, derive every namespace from it, and refuse an overlay that states a derived name. Two
instances must share nothing, with locations compared by real path and containment and secrets
servers by the host the resolver reaches. Named instances are disjoint by construction; the default
instance is checked at load. The tests fail under five broken derivations and six loader mutations,
and two bypasses found in a security review are fixed. The check holds at load time only.
\end{abstract}

\section{Introduction}

A knowledge store built on the mind-palace framework keeps its state where its configuration says:
a data directory holding its stores, queues and lock; a file of parameters written by a self-tuning
loop; an outbound directory that a sync service mirrors off the host; secrets in the platform
keychain under fixed item names; two services registered with the service manager under fixed
labels; a remote backup repository; and, optionally, a secrets server. The configuration is a chain
of TOML files merged in order: committed defaults, then the tuned parameters, then a hand-written
overlay that wins over both.

Before this change nothing in that chain said which instance it described. The overlay's file name
was the only mark of identity, so a copied checkout started on another host, or a second
configuration started beside the first, resolved each of those names to the same place and became
the same individual. It would append to the same history, read the same corpus, hold the same keys
and write the same backup repository. A self-tuning loop running in the second instance would also
write the first instance's parameter file, which the first reads at its next start.

This paper describes how instance identity was made explicit and mechanical. The contributions are:
\begin{enumerate}
\item \emph{An identity.} Every named overlay declares an instance id from a grammar that makes it
  one path segment and one dotted component and keeps distinct ids distinct on a case-insensitive
  file system. One id is reserved to the default chain, which keeps every name it had
  (\cref{sec:model}).
\item \emph{A derivation.} Every namespace a named instance holds is a function of its id: data
  root, lock, spool, parameter file, outbound lane, secret names, service labels and the last
  segment of its backup repository. An overlay that states a derived name is refused, not
  overridden (\cref{sec:concepts}).
\item \emph{An invariant and its check.} For instances $i \neq j$, $\shared(N(i), N(j)) =
  \emptyset$, where locations compare by containment of their real paths and a secrets server
  compares by the host the resolver would reach. Between named instances the derived names are
  disjoint by construction; against the default instance the loader evaluates the invariant at
  load (\cref{sec:model,sec:impl}).
\item \emph{Evidence that the checks fail on wrong code:} a Hypothesis property through the real
  loader with five broken derivations as controls, six mutations of the loader each failing a named
  subset of 183 tests, a differential comparison of the IPv4 parser with the platform's
  \texttt{inet\_aton} on $10^6$ strings, and the cost of the check (\cref{sec:eval}).
\item \emph{Two bypasses} found by a security review of the branch, their fixes, and the limits
  that remain (\cref{sec:audit,sec:limits}).
\end{enumerate}

\section{Background and related work}

\subsection{The configuration chain}
Let $\oplus$ be the deep merge of two TOML documents in which the right operand wins on every leaf.
The default chain is $C_d = D \oplus L_d \oplus O_d$: the committed defaults $D$, the parameter file
$L_d$ written by the self-tuning loop, and the hand-written overlay $O_d$, applied last so that a
human setting always wins over a tuned one. The default chain is one instance, the one that existed
before instances had ids. We call it the default instance and write its reserved id as $d$. Paths
in the chain are relative to the repository root $P$ unless absolute.

A second instance is selected at run time by an environment variable, \env, that names one overlay
file $O_i$. That overlay replaces both $L_d$ and $O_d$: a named instance is a different individual,
so it inherits neither the default instance's hand settings nor the parameters its loop tuned. A
variable that names no file is refused rather than falling back to the committed defaults, since a
silent fallback would bring the instance up with the wrong embedding model and similarity
threshold.

\subsection{Related work}

\paragraph{Protection principles.} Saltzer and Schroeder~\cite{saltzer1975} list, among their design
principles, least common mechanism (minimize the mechanism common to more than one user) and
fail-safe defaults (base access decisions on permission rather than exclusion). A named instance
holds no name in common with another, and it reads no corpus it did not declare, which is a
fail-safe default in their sense. The secrets-server comparison is not one: it refuses a server
that canonicalizes to the default instance's and accepts every other, so its soundness rests on the
canonicalization. The first version recognized the loopback host by a list of spellings, and the
security review found spellings missing from the list (\cref{sec:audit}).

\paragraph{Operating-system namespaces.} Plan~9 gives each process its own name space, assembled by
mounts and binds~\cite{pike1993names}. FreeBSD jails partition a host into environments with their
own file system root, processes and addresses~\cite{kamp2000jails}. For Linux, Biederman proposes
that each class of global identifier, such as process ids, System~V IPC keys and user ids, become a
namespace with multiple instances~\cite{biederman2006}. These mechanisms are enforced by the
kernel, so a confined
process cannot name what lies outside its view. The mechanism here is weaker and sits higher. It
decides which names an instance's configuration resolves to, within one user account, and it cannot
stop a process from opening a path it was not given. The two compose: the derived root is one
directory per instance, which is what a container profile mounts as the instance's only volume.

\paragraph{Time of check and time of use.} Bishop and Dilger~\cite{bishop1996} study flaws in which
the binding of a file name changes between a check and the use that relies on it. The check here
runs once, when the configuration is loaded, and a symbolic link planted afterwards changes what a
checked name refers to. We state this as a limit (\cref{sec:limits}).

\paragraph{Parser differentials.} When the component that validates an address parses it
differently from the component that connects to it, the validation can be bypassed;
Tsai~\cite{tsai2017} demonstrated this across the URL parsers and HTTP clients of many languages.
The grammar that matters here is the URI authority, in which an optional userinfo ends with
\texttt{@} before the host~\cite{rfc3986} (the comparison takes everything up to the last
\texttt{@} as userinfo, so no spelling hides a host behind one); the rule that \texttt{localhost}
and its subdomains name the loopback
address~\cite{rfc6761}; IPv4-mapped IPv6 addresses~\cite{rfc4291}; and the permissive IPv4 forms the
C library's \texttt{inet\_aton} accepts. The canonicalizer of \cref{alg:host} reads a host the way
the resolver will, not the way a regular expression does.

\paragraph{Property-based and mutation testing.} QuickCheck~\cite{claessen2000} states properties as
executable predicates checked on generated inputs; Hypothesis~\cite{maciver2019} brings the approach
to Python with shrinking and composable strategies. A property is evidence only to the extent that it
fails on wrong code. Mutation analysis~\cite{demillo1978} measures a test suite by the deliberate
faults it detects. We use both: the property's controls are hand-written mutants of the derivation,
and we report six further mutations of the loader (\cref{sec:eval}).

\section{Model}\label{sec:model}

No Lean declaration checks the statements in this section. They are checked by the tests and
measurements of \cref{sec:eval}.

\begin{definition}[Instance id]
An instance id is a string matching \texttt{[a-z][a-z0-9-]\{0,31\}}: a lowercase letter followed by
up to 31 lowercase letters, digits or hyphens. The id $d$ is reserved to the default chain. A named
overlay must declare an id other than $d$, in an \texttt{[instance]} table that holds no other key.
The default chain, and a file loaded by explicit path, may restate $d$ and nothing else.
\end{definition}

\begin{definition}[Derivation]
For an id $i$ let $\rho(i) = P$ if $i = d$, and $\rho(i) = P/\texttt{instances}/i$ otherwise. The
secret name of $n$ for $i$ is $\sigma(i, n) = n$ if $i = d$ and the string $i.n$ otherwise. The
service label of role $r$ is $\lambda(i, r) = \texttt{com.mind-palace.}r$ if $i = d$ and
$\texttt{com.mind-palace.}i.r$ otherwise. A relative data path $x$ resolves to $\rho(i)/x$. A named
instance's chain is $C_i = D \oplus L_i \oplus O_i$ with $L_i = \rho(i)/\texttt{levers.toml}$.
\end{definition}

\begin{definition}[Namespaces]
The namespaces of instance $i$ are the tuple
\[
  N(i) = (W_i, R_i, Q_i, S_i, \Lambda_i, b_i, A_i, v_i)
\]
read off the configuration the loader builds: $W_i$, the locations it writes (14 data
paths, the supervisor lock, the parameter file and the outbound lane); $R_i$, the corpus locations
it reads (a notes vault and a transcripts directory); $Q_i$, its dialogue lanes, readers over the
histories of two local chat applications; $S_i$, its secret names (five keychain items, a cloud
credential profile and the secrets-server token item); $\Lambda_i$, its service labels; $b_i$, its
backup repository or $\varepsilon$; $A_i$, the object-store prefixes of its request and result
exchange; and $v_i$, its secrets server as \texttt{host:port} or $\varepsilon$.
\end{definition}

\begin{definition}[Location and nearness]
For a path $p$, $\loc(p)$ is $p$ made absolute, with every symbolic link on it resolved and its
case folded; a path that does not exist resolves through its deepest existing ancestor with the
rest appended. For strings $x, y$ and a set of separators $\Sigma$, write $x \sqsubseteq_\Sigma y$
when $x = y$ or $x$ begins with $y\,s$ for some $s \in \Sigma$. The strings are near, $x
\near_\Sigma y$, when $x \sqsubseteq_\Sigma y$ or $y \sqsubseteq_\Sigma x$.
\end{definition}

\begin{definition}[Shared]
$\shared(N(i), N(j))$ is the set made of: every pair $(p, q) \in (W_i \cup R_i) \times (W_j \cup
R_j)$ with $\loc(p) \near_{\{/\}} \loc(q)$; every element of $Q_i \cap Q_j$, $S_i \cap S_j$ and
$\Lambda_i \cap \Lambda_j$; the pair $(b_i, b_j)$ if both are set and their case-folded texts,
without a trailing slash, are near under $\{/, :\}$; every pair in $A_i \times A_j$ near under
$\{/\}$, compared the same way; and $v_i$ if $v_i \neq \varepsilon$ and $v_i = v_j$, hosts canonicalized by
\cref{alg:host}.
\end{definition}

\begin{definition}[Disjointness invariant]\label{def:invariant}
For instances $i \neq j$, $\shared(N(i), N(j)) = \emptyset$.
\end{definition}

Containment, not equality, is the relation that matters for locations: a store inside another
instance's data directory rides that instance's backup and falls to its reset. The comparison is
also deliberately reflexive. $\shared(N(i), N(i))$ is non-empty whenever $W_i \neq \emptyset$, and
the property test asserts this, so a comparison that always returns nothing cannot pass it.

\begin{lemma}[The id grammar]\label{lem:grammar}
Let $i \neq j$ be ids. (a) Neither contains \texttt{/} or \texttt{.}, and case folding fixes both.
(b) $P/\texttt{instances}/i$ and $P/\texttt{instances}/j$ are not near under $\{/\}$. (c) For named
$i$ and $j$, $\sigma(i, n) \neq \sigma(j, m)$ and $\lambda(i, r) \neq \lambda(j, r')$ for all names
and roles.
\end{lemma}
\begin{proof}[Proof sketch]
(a) is the grammar. (b) The two paths share the prefix $P/\texttt{instances}/$ and differ in their
last segment. One is near the other only if one segment equals the other, or begins with it
followed by \texttt{/}, and (a) excludes both. (c) A named secret name splits at its first dot into
the id and the rest, because the id contains no dot, so equal names have equal ids. A named label's
third dot-separated component is the id, because \texttt{com} and \texttt{mind-palace} contain no
dot.
\end{proof}

\begin{proposition}[Named instances are disjoint by construction]\label{prop:named}
Let $i \neq j$ be named ids whose overlays declare no corpus source, backup repository, object-store
bucket or secrets server, and suppose no symbolic link exists at or below $P/\texttt{instances}$.
Then $\shared(N(i), N(j)) = \emptyset$.
\end{proposition}
\begin{proof}[Proof sketch]
A load that returns has refused every absolute data path and every relative one containing
\texttt{..}, so each data path of $i$ lies under $\rho(i)$. So do the lock, the spool, the
parameter file, the outbound lane (derived, or moved by the overlay within $\rho(i)$) and the two
undeclared sources, $\rho(i)/\texttt{vault}$ and $\rho(i)/\texttt{transcripts}$. With no links
below $P/\texttt{instances}$, $\loc$ preserves the prefix $\rho(i)$, so by
Lemma~\ref{lem:grammar}(b) no location of $i$ is near a location of $j$. Secret names and labels differ
by Lemma~\ref{lem:grammar}(c), the dialogue lanes of a named instance are empty, and $b$, $A$ and $v$ are
empty by hypothesis.
\end{proof}

\begin{proposition}[The default instance is checked at load]\label{prop:default}
If the loader returns a configuration for a named instance $i$, then $\shared(N(i), N(d)) =
\emptyset$, where $N(d)$ is read off the default chain present on the host at the moment of
loading.
\end{proposition}
\begin{proof}
The last step of the load computes $N(d)$ from the default chain and refuses when the set is not
empty (\cref{alg:load}, line~\ref{line:refuse}).
\end{proof}

The proof of Proposition~\ref{prop:default} is one line because its content is in what $N(d)$ covers. The
default instance's configuration is written by hand and may place a store, a source or a repository
anywhere, so its names cannot be derived and disjointness from it cannot follow from
Proposition~\ref{prop:named}. The check covers what Proposition~\ref{prop:named} assumes away
(declared sources,
repositories, prefixes and servers, and symbolic links) by comparing against the default instance's
actual names. It also covers the default instance's own secret names: the committed names contain no
dot, so none equals a derived $i.n$, but a hand overlay that renamed one to such a form would be
caught here. Named overlays are not compared with each other at load; \cref{sec:pairs,sec:limits}
state what that leaves open.

\paragraph{The secrets-server host.} A server is compared as $\mathrm{canon}(h){:}\mathit{port}$,
where a missing port is the scheme's and $\mathrm{canon}$ (\cref{alg:host}) maps every spelling of
the local machine to one token, \texttt{loopback}: the names \texttt{localhost} and
\texttt{*.localhost}~\cite{rfc6761}, every loopback or unspecified IPv4 or IPv6 address,
IPv4-mapped addresses read as their IPv4 address~\cite{rfc4291}, and every IPv4 form
\texttt{inet\_aton} accepts. The unspecified addresses count because connecting to them reaches the
local machine. Another address maps to its canonical text, so two spellings of one remote server
also compare equal; a name that is not an address stays as written, case folded and without one
trailing dot. The intended property is agreement with the resolver: whenever the platform would
connect two spellings to one IPv4 address, they canonicalize to one string. We test this property
rather than prove it (\cref{sec:aton}), because its reference is the platform's C library.

\section{Concepts and intuition}\label{sec:concepts}

\paragraph{One string, many names.} Every name an instance holds is computed from one string.
Whether two instances collide then reduces to whether two ids are equal, which the grammar makes a
plain string comparison: no case variants, no path separators, and no dots that would let one
derived name parse as another.

\paragraph{Derived, not declared.} A named overlay states none of the names its id derives. One
that does is refused, even when it states the derived value, and is not overridden: a value the
loader silently replaces is a setting that looks applied and never took. A declared name would also
be a second source of truth that could collide with another instance's; a derived one cannot. The
backup repository is the one declared location among the written ones, because its place is the
operator's choice; only its last segment is constrained to be the id.

\paragraph{Sources are declared or absent.} Locations an instance reads are treated differently
from those it writes. A named instance may read a corpus anywhere except at, inside or around a
location of the default instance. With nothing declared it reads nothing, because the committed
defaults name the default instance's corpus. An agent enabled without a declared source is refused,
and two named instances may deliberately read one corpus.

\paragraph{Compare where things are, not what they are called.} Two names for one directory are one
directory, and a directory inside another is written whenever the outer one is backed up or reset.
The comparison therefore resolves symbolic links and treats containment as sharing, and it compares
a secrets server by the host the resolver would reach. Both rules were added after the security
review showed that the lexical versions could be defeated (\cref{sec:audit}).

\paragraph{The default instance is read only to refuse.} A named load builds the default instance's
configuration for one purpose, computing $N(d)$, and no value from it reaches the named
configuration. The default instance's names, which predate ids, are exactly what they were.

\paragraph{Fail closed.} Every guard raises; none warns. Tools that reach the default instance's
state by a fixed path rather than through the configuration refuse to run under a named overlay
instead of learning about instances. Making one of them instance-aware later is a change to read
the configuration.

\section{Implementation}\label{sec:impl}

The loader is one module. \Cref{alg:load} is the path a named overlay takes. The default chain takes
an earlier branch, unchanged except that it records the id $d$; a file passed by explicit path is
read verbatim and is also the default instance.

\begin{algorithm}
\caption{Loading a named overlay.}\label{alg:load}
\begin{algorithmic}[1]
\Require $O$, the file named by \env
\Ensure a configuration for the instance that $O$ declares, or a refusal
\State $X \gets \Call{Read}{O}$; \quad $i \gets$ the id in the \texttt{[instance]} table of $X$
\If{$i$ is absent or malformed, or $i = d$} \textbf{refuse} \EndIf
\State $T \gets \Call{Read}{D}$; \quad $K \gets$ the committed secret names in $T$
\State $T \gets T \oplus \Call{Read}{\rho(i)/\texttt{levers.toml}} \oplus X$
\State point each undeclared source at $\rho(i)/\mathit{name}$ and disable it; refuse an enabled
  agent without a source
\State give an undeclared secrets server the empty address; refuse secrets enabled without one
\State refuse if $X$ names a key of $K$, or $T$ holds for it anything but the committed value
  \label{line:secret}
\State refuse if a data path or the outbound lane is absolute, starts at \texttt{\textasciitilde},
  or contains \texttt{..}
\State refuse if a backup repository is set and its last segment is not $i$
\State $C \gets \Call{Build}{T, i}$ \Comment{paths under $\rho(i)$, names $\sigma(i, \cdot)$}
\State $C_d \gets \Call{Build}{D \oplus L_d \oplus O_d,\ d}$ \Comment{read only for $N(d)$}
\If{$\Call{Shared}{N(C), N(C_d)} \neq \emptyset$} \label{line:refuse}
  \textbf{refuse}, listing each shared name
\EndIf
\State \Return $C$
\end{algorithmic}
\end{algorithm}

Line~\ref{line:secret} refuses a secret name when the overlay names the key at all, even with the derived value,
and also when the merged chain holds anything but the committed default for it; the second clause
catches a parameter file that sets one. The comparison $\shared$ is the definition in
\cref{sec:model} written as nested loops, so every location of one instance is resolved and
compared with every location of the other. \Cref{alg:host} is the canonical host. Its IPv4 parser
is written out by hand, because a static import check forbids networking modules in the core
package and \texttt{socket.inet\_aton} is out of reach; it follows the BSD rules.

\begin{algorithm}[t]
\caption{The host a secrets-server address reaches.}\label{alg:host}
\begin{algorithmic}[1]
\Function{Canon}{$h$} \Comment{$h$: the authority after its last \texttt{@}, unbracketed}
\State decode percent escapes; fold case; drop one trailing dot
\If{$h$ is empty, \texttt{localhost}, or ends in \texttt{.localhost}} \Return \texttt{loopback}
\EndIf
\If{$h$ without an IPv6 zone parses as an IP address} $a \gets$ that address
\ElsIf{$\Call{InetAton}{h} \neq \bot$} $a \gets$ that IPv4 address
\Else{} \Return $h$
\EndIf
\If{$a$ is IPv4-mapped} $a \gets$ its IPv4 address \EndIf
\If{$a$ is loopback or unspecified} \Return \texttt{loopback} \EndIf
\State \Return the canonical text of $a$, bracketed if IPv6
\EndFunction
\Statex
\Function{InetAton}{$t$} \Comment{the BSD C library's reading}
\State read up to four dot-separated parts, each starting with a digit: decimal, octal after
  \texttt{0}, hexadecimal after \texttt{0x}; the 64-bit accumulator wraps
\State fail on an 8 or 9 in an octal part, a leading part above 255, a last part with no digit, a
  character after the address other than whitespace, or, with $k > 1$ parts, a last part that
  does not fit the $5 - k$ bytes left to it
\State \Return the leading parts as the high bytes over the last; one part keeps its low 32 bits
\EndFunction
\end{algorithmic}
\end{algorithm}

\paragraph{Run-time guards.} Two kinds of code reach an instance's names outside the loader. The
self-tuning loop writes the parameter file; it now receives the path from the configuration, and a
loop built for a named instance with any other path refuses each step before writing. Tools that
reach the default instance's state by a fixed path (its data directory, its services, its archive,
its secrets server, and the default output directory of the report writers) call a one-line guard
as their first statement, which refuses while \env{} is set. Shell scripts carry the same guard
inline, as do two Python entry points that never import the core. Among the consumers of the
derived names, the service launcher takes its job files, labels and lock from them, and the report
writer places a named instance's reports in its own lane; the stack, backup and secret-loading
scripts hold the default instance's labels and items and refuse a named instance.

\section{Evaluation}\label{sec:eval}

\subsection{Setup}
All runs used macOS 27.0.1 on an Apple M2 Pro with 32\,GiB of memory, Python 3.15.2, pytest 9.1.1
and Hypothesis 6.155.7, at the branch head \texttt{babe4de} over the base \texttt{f65f73f}.
Mutation runs used an export of the branch head (\texttt{git archive}) with one edit applied at a
time, and ran all 183 instance tests for each edit. The scripts for the mutations, the differential
comparison and the timing are short and are described with each result, so each run can be
repeated from the description.

\begin{table}[ht]
\centering\small
\caption{The instance tests at the branch head: 183 pass in 9.3\,s.}\label{tab:tests}
\begin{tabular}{@{}llr@{}}
\toprule
Kind & What it checks & Tests \\
\midrule
unit & overlay selection, ids, refusals, derived names, the pinned default, servers & 130 \\
integration & links, parameter files, consumers of the names, tool guards & 41 \\
property & the invariant, its controls, repositories, IPv4 spellings & 12 \\
\midrule
total & & 183 \\
\bottomrule
\end{tabular}
\end{table}

\subsection{The property and its controls}
The property draws lists of one to three distinct ids from the id grammar, excluding $d$. For each
id it writes an overlay declaring only the id and a repository ending in it, builds the
configuration through the real loader, adds the default instance, and asserts that every instance
shares something with itself and that every pair shares nothing. It runs 40 examples per execution
with no example database, so each run draws new ids. A second property, over 60 generated ids,
checks that the default instance's repository, named outright or with a segment below it, is
refused for every id, while a location of the instance's own whose text extends the default
repository's without a separator loads.

Each control replaces one function of the derivation with a broken version, switches off the
load-time refusal so that the property sees the broken names, and asserts that the property fails
with the expected shared component (\cref{tab:controls}). With the refusal switched on, four of the
five are refused by the loader itself. The fifth, which gives every named instance the same root,
prefix and label, is caught only by the property: the loader compares a named instance with the
default instance and never with another named one (\cref{sec:pairs}).

\begin{table}[ht]
\centering\small
\caption{Broken derivations used as controls. Each fails the property.}\label{tab:controls}
\begin{tabular}{@{}lll@{}}
\toprule
Broken derivation & Reported as shared & Refused at load \\
\midrule
id ignored (one root, prefix and label for all) & a location & no (a named pair) \\
named instance given the default's names & a location & yes \\
undeclared vault becomes the default's & the default's vault & yes \\
undeclared outbound lane becomes the default's & the default's lane & yes \\
no server named: the committed address & server \texttt{loopback:8200} & yes \\
\bottomrule
\end{tabular}
\end{table}

\subsection{Mutations of the loader}
\Cref{tab:mutations} lists six single edits to the loader and the tests each one fails. The first
two restore the code as it was before the security review. Every mutation is detected, and each
fails only tests written for the behaviour it removes, so the tests are specific as well as
sensitive.

\begin{table}[ht]
\centering\small
\caption{Single mutations of the loader, each run against all 183 tests.}\label{tab:mutations}
\begin{tabular}{@{}lrl@{}}
\toprule
Mutation & Fail & Failing tests \\
\midrule
M1 lexical paths, not real ones & 4 & the four link cases \\
M2 the earlier server comparison & 31 & 14 spellings, 2 tests each; 3 remote hosts \\
M3 load-time refusal removed & 42 & links 4, controls 4, repositories 3, sources 5, server 26 \\
M4 equality in place of containment & 6 & nested repositories 3, sources 2, a link 1 \\
M5 id check removed from the id type & 4 & malformed ids built without the loader \\
M6 one number not cut to 32 bits & 3 & the spelling past 32 bits, in 3 tests \\
\bottomrule
\end{tabular}
\end{table}

\subsection{The security review}\label{sec:audit}
A security review of the branch, made after the derivation and the check were in place, reported
two ways to make the check pass while the instance still reached the default instance's state.

\paragraph{Symbolic links.} The comparison was lexical: absolute, normalized and case folded. With
$\rho(i)$ linked to $P$, the instance's data directory read as $P/\texttt{instances}/i/\texttt{data}$
while it was $P/\texttt{data}$, and nothing was shared. Linking $P/\texttt{instances}$, or the
instance's own data directory, into $P/\texttt{data}$ did the same, and so did a declared source that
was a link to the default instance's vault. The fix compares real paths. Under M1 exactly these four
tests fail; a control shows that a root linked to a place of the instance's own still loads.

\paragraph{Spellings of the loopback address.} The comparison parsed the address with a regular
expression and recognized the local host by a fixed set (the empty host, \texttt{localhost},
\texttt{::1} and its long form, \texttt{0.0.0.0}, \texttt{::}) and the prefix \texttt{127.}.
Userinfo stayed part of the host, and no numeric form other than a dotted \texttt{127.} address was
recognized. \Cref{tab:spellings} lists the 14 spellings of the local server the earlier comparison
did not fold. Each appears in two tests, one for the canonical text and one for the refusal at load,
which gives 28 of the 31 failures under M2. The other three are controls for remote hosts whose
canonical text the earlier code did not produce: \texttt{010.0.0.5}, which the resolver reads as
octal \texttt{8.0.0.5}; \texttt{[::2]}; and \texttt{[::ffff:10.0.0.1]}, which is
\texttt{10.0.0.1}.

\begin{table}[ht]
\centering\small
\caption{Spellings of the local secrets server that the earlier comparison did not fold to
\texttt{loopback}; all fold now.}\label{tab:spellings}
\begin{tabular}{@{}llr@{}}
\toprule
Class & Spellings tested (host part) & Count \\
\midrule
userinfo before the host & \texttt{u@127.0.0.1}, \texttt{u:p@localhost}, \texttt{a@b@127.0.0.1} & 3 \\
one number & \texttt{2130706433}, \texttt{0x7f000001}, \texttt{017700000001} & 3 \\
one number past 32 bits & \texttt{6425673729} & 1 \\
hexadecimal dotted form & \texttt{0x7f.1} & 1 \\
the unspecified address as one number & \texttt{0} & 1 \\
trailing dot & \texttt{localhost.} & 1 \\
subdomain of \texttt{localhost} & \texttt{vault.localhost} & 1 \\
percent-encoded name & \texttt{\%6c\%6f\%63\%61\%6c\%68\%6f\%73\%74} & 1 \\
IPv4-mapped IPv6 & \texttt{[::ffff:127.0.0.1]}, \texttt{[::ffff:7f00:1]} & 2 \\
\midrule
total & & 14 \\
\bottomrule
\end{tabular}
\end{table}

\subsection{The IPv4 parser against the platform}\label{sec:aton}
A Hypothesis property writes each of 300 random 32-bit addresses per run nine ways (dotted decimal,
one decimal number, one hexadecimal number, one octal number, the two- and three-part short forms,
dotted hexadecimal, dotted octal, and the decimal number plus $2^{32}$) and checks that every
spelling canonicalizes to the address, or to \texttt{loopback} for $127.0.0.0/8$ and $0.0.0.0$. On
every spelling the platform accepts, it also compares the parse with \texttt{socket.inet\_aton}.
BSD's \texttt{inet\_aton}, which macOS uses, reads a single number past 32 bits as its low bits;
glibc's refuses it. The parser keeps the BSD reading, the wider one, so a spelling that either
library resolves to the local machine folds to \texttt{loopback}; on Linux the ninth spelling is
compared with the expected text only.

Beyond the property, we compared the parser with the platform's \texttt{inet\_aton} on $10^6$
random strings of length 1 to 16 over digits, the hexadecimal letters in both cases, \texttt{x},
\texttt{X}, the dot, space and tab, drawn with a fixed seed. Both accepted 214{,}495 strings with
equal values, both refused 785{,}505, and they disagreed on none. The run took 3.3\,s.

\subsection{The default instance is unchanged}
The unit suite pins the default instance on a fixture that has the sections and keys of a real
default overlay with placeholder values, holding each written path, secret name and the repository
to values recomputed from the TOML; a control loads the same overlay as a named instance and finds
every namespace different. Independently, we loaded that fixture as the default chain with the base
commit's loader and with the branch's, serialized each configuration with the new id field removed
and the repository root replaced by a token, and compared the two: 22 sections and 156 leaf values,
identical.

\subsection{Cost}
\Cref{tab:cost} gives the cost of a load. The named load builds two configurations and compares 19
locations of each instance pairwise, resolving each location again for every pair: 722 real-path
resolutions. Under a profiler (50 loads) the check is 87\% of a named load, and the real-path
resolutions alone are 66\%. A load happens once per process start, so the cost is paid once.

\begin{table}[ht]
\centering\small
\caption{Wall time of one load: medians of 200 loads after a warm-up, in three
repetitions.}\label{tab:cost}
\begin{tabular}{@{}lrl@{}}
\toprule
Load & Median & Range of the three medians \\
\midrule
default chain (committed defaults, no overlay present) & 0.83\,ms & 0.80 to 0.85\,ms \\
named instance (the committed benchmark overlay) & 10.7\,ms & 10.62 to 10.82\,ms \\
\bottomrule
\end{tabular}
\end{table}

\subsection{Two named instances}\label{sec:pairs}
The loader compares a named instance with the default instance only. We checked what that leaves
open on a throwaway repository layout. Two named instances whose repositories nest
(\texttt{\dots/alpha} and \texttt{\dots/alpha/beta}) both load, since each ends in its own id, and
$\shared$ between them reports the repository pair. A named root that is a link to another named
root ($P/\texttt{instances/gamma}$ to $P/\texttt{instances/alpha}$) also loads, and $\shared$
between the two reports 49 shared locations. In both cases the comparison detects the overlap; the
loader does not ask it to. Proposition~\ref{prop:named} excludes both cases by hypothesis.

\section{Limitations}\label{sec:limits}

\paragraph{The check holds at load.} A symbolic link planted after the load changes what a checked
name refers to, and a path that did not exist at load was resolved through its deepest existing
ancestor, so a link created later at a missing component is not seen. This is the time-of-check,
time-of-use pattern~\cite{bishop1996}, and the check does not claim to close it.

\paragraph{Host names are not resolved.} The core package resolves no names, so a DNS name that
resolves to the local machine passes the server comparison. The named instance's token item is its
own, so such a name would reach the default instance's server without the default instance's
token.

\paragraph{Named instances are not compared with each other.} Proposition~\ref{prop:named} covers
their derived
names. Their declared names are compared with the default instance's only: two named instances may
read one corpus by design, and may also name nested repositories, overlapping object-store prefixes
or one secrets server, and a link between two named roots goes unseen (\cref{sec:pairs}). The
repository rule constrains only the last segment, which keeps two repositories distinct but not
apart.

\paragraph{Locking is per host.} The supervisor lock is an advisory file lock in the instance's data
directory. It excludes a second supervisor of one instance on one host, not on two hosts that share
storage.

\paragraph{The comparison needs the default overlay.} In a container image without the default
overlay, the default chain is the committed defaults alone, so the check cannot see the names that
overlay sets by hand, such as its backup repository. The named instance's derived password still
keeps that repository closed, since the backup tool neither opens a repository with another
password nor initializes over it.

\paragraph{Conservative choices.} Case folding matches the macOS default file system; on a
case-sensitive one it can report two distinct directories as shared and refuse a load that would be
safe, never the reverse. The IPv4 parser's BSD reading can likewise fold to \texttt{loopback} a
spelling that glibc would refuse.

\paragraph{What a named instance cannot yet do.} No configuration can declare a dialogue lane for a
named instance, so it has none. The tools that reach the default instance's state by path refuse
rather than serve a named instance.

\section{Future work}
Comparing a named instance with every overlay known on the host, not only with the default chain,
would close the gap of \cref{sec:pairs}: $\shared$ already detects both demonstrated overlaps, and
the cost grows with the number of pairs. A registry of instances, planned for the key fabric of a
later step, would serve the same purpose across hosts and could hold a lease in place of the
per-host lock. Resolving the secrets server's name at connection time, in the networked edge rather
than the core, and pinning the client to the address that was checked would close the DNS gap.
Opening state through directory descriptors that refuse to follow links, or checking each store
again when it is opened, would narrow the window after load. Resolving each location once per load
would reduce the 722 real-path resolutions to 38. Finally, Lemma~\ref{lem:grammar} and
Proposition~\ref{prop:named} are small enough to state and check in Lean over an abstract model of
paths and
names.

\section{Conclusion}
An instance id, required of every named overlay and reserved for the default chain, turns identity
from a file name into a value that every namespace is computed from. Derived names make two named
instances disjoint by construction, and a load-time comparison against the default instance's actual
configuration covers what derivation cannot: hand-placed locations, declared sources and servers,
and links. The comparison was made to read locations by their real paths and hosts by the
resolver's rules after a review showed the lexical versions could be bypassed, and every check
reported here fails on code that lacks it. The guarantees end at load time, at the host boundary,
and at the boundary between two named instances, and the paper states each of those limits.

\bibliographystyle{plainnat}
\bibliography{refs}
\end{document}
